\section{Comprehensive Scaling Blueprint to Real-World Datasets and Architectures}
\label{app:scaling_blueprint}

In this section, we present a complete, step-by-step mathematical and procedural blueprint for scaling \textbf{PAC-Bayes Merge} beyond the simulated representation sandbox to physical computer vision backbones and autoregressive Large Language Models (LLMs). This demonstrates that our theoretical framework is general-purpose, cross-modality, and directly applicable to large-scale deep learning models.

\subsection{Physical Computer Vision backbones (e.g., ResNets, Vision Transformers)}
For physical computer vision tasks, rather than operating on low-dimensional synthetic representations, PAC-Bayes Merge operates directly on the parameter weight tensors of pre-trained and fine-tuned deep neural networks (such as ResNet-18, ResNet-50, or ViT-B/16). 

The step-by-step implementation procedure is as follows:
\begin{enumerate}
    \item \textbf{Expert Model Fine-Tuning:} Start from a shared pre-trained Vision Transformer $W_0$ (e.g., \texttt{vit-base-patch16-224}). Independently fine-tune $K$ copies of the network on $K$ physical vision datasets (e.g., CIFAR-10, CIFAR-100, SVHN, and Flowers-102) to obtain task-specific expert parameter weights $W_k$ for $k \in \{0, \dots, K-1\}$.
    \item \textbf{Layer Grouping and Task Vector Extraction:} Group the network parameters into $L$ discrete blocks (e.g., the 12 Transformer blocks of a ViT-Base). For each layer group $l \in \{0, \dots, L-1\}$ and task $k$, extract the weight-space task vector:
    \begin{equation}
    V_k^{(l)} = W_k^{(l)} - W_0^{(l)}
    \end{equation}
    where $W_k^{(l)}$ and $W_0^{(l)}$ represent the flattened weight tensors of block $l$.
    \item \textbf{Trajectory Expansion:} Initialize the learnable trajectory coordinates $\Theta \in \mathbb{R}^{K \times (d+1)}$ with the uniform target $\theta_{k,0} = \ln(1/(K-1))$ and $\theta_{k,j \ge 1} = 0.0$. At each training step, compute the ensembling coefficients $\alpha_k(l; \theta_k) = \sigma\left(\sum_{j=0}^d \theta_{k,j} (l/(L-1))^j \right)$ across all layers.
    \item \textbf{Monte Carlo Weight Assembly:} Draw $S$ independent Monte Carlo samples of the trajectory parameters $\tilde{\Theta}_s \sim \mathcal{N}(\Theta, \sigma^2 I)$ and assemble the corresponding physical merged weights:
    \begin{equation}
    W_{\text{merged}, s}^{(l)} = W_0^{(l)} + \sum_{k=0}^{K-1} \alpha_k(l; \tilde{\theta}_{k, s}) V_k^{(l)}
    \end{equation}
    \item \textbf{Expected Risk Minimization:} Pass a tiny few-shot calibration dataset (e.g., $M = 10$ real pixel images per class, resized and normalized to $224 \times 224$) through the network with the assembled merged weights. Backpropagate the average cross-entropy loss over the $S$ samples to update $\Theta$:
    \begin{equation}
    \mathcal{L}_{\text{total}}(\Theta) = \frac{1}{S \cdot K \cdot M} \sum_{s=1}^S \sum_{k=0}^{K-1} \sum_{i=1}^M \ell_{\text{ce}}\left(f\left(x_i^{(k)}; W_{\text{merged}, s}\right), y_i^{(k)}\right) + \lambda_{\text{PAC}} \|\Theta - \Theta_{\text{uniform}}\|_2^2
    \end{equation}
    \item \textbf{Static Compilation:} After $T$ optimization epochs (e.g., $T=100$), compute the optimal mean coefficients $\alpha_k^*(l) = \sigma(p_k(l/(L-1); \theta_k^*))$ and statically compile the final model weights $W_{\text{final}}^{(l)} = W_0^{(l)} + \sum_k \alpha_k^*(l) V_k^{(l)}$ for zero-latency, zero-overhead deployment on real physical test images.
\end{enumerate}

\subsection{Autoregressive Large Language Models (LLMs) with LoRA Merging}
In the Large Language Model community, full-parameter model merging is often prohibitive due to the immense scale of the weights (e.g., LLaMA-3 8B, Mistral 7B). However, post-hoc merging is extremely popular for combining specialized Low-Rank Adaptation (LoRA) adapters~\cite{hu2021lora} fine-tuned from a shared pre-trained base LLM.

PAC-Bayes Merge is perfectly compatible with LoRA adapter merging. For each layer $l$, let $W_0^{(l)} \in \mathbb{R}^{D_{\text{out}} \times D_{\text{in}}}$ represent the pre-trained query, key, value, or projection weight matrix of the autoregressive Transformer. Suppose we have $K$ task-specific LoRA adapters (e.g., one fine-tuned for Python coding, one for medical Q\&A, one for mathematical reasoning, and one for French translation), each parameterized by low-rank matrices $B_k^{(l)} \in \mathbb{R}^{D_{\text{out}} \times r}$ and $A_k^{(l)} \in \mathbb{R}^{r \times D_{\text{in}}}$ where $r \ll \min(D_{\text{out}}, D_{\text{in}})$.

We define two mathematically elegant LoRA merging strategies under our trajectory-regularized framework:
\begin{enumerate}
    \item \textbf{Direct Parameter-Space LoRA Merging:}
    We parameterize the ensembling coefficients $\alpha_k(l; \theta_k)$ as before. Since LoRA task vectors are represented as product matrices, the merged adapter weights are computed as:
    \begin{equation}
    W_{\text{merged}}^{(l)}(\Theta) = W_0^{(l)} + \sum_{k=0}^{K-1} \alpha_k(l; \theta_k) B_k^{(l)} A_k^{(l)}
    \end{equation}
    Alternatively, to maintain the low-rank factorized format during inference, we can merge the adapter matrices directly:
    \begin{equation}
    B_{\text{merged}}^{(l)}(\Theta) = \sum_{k=0}^{K-1} \sqrt{\alpha_k(l; \theta_k)} B_k^{(l)}, \quad A_{\text{merged}}^{(l)}(\Theta) = \sum_{k=0}^{K-1} \sqrt{\alpha_k(l; \theta_k)} A_k^{(l)}
    \end{equation}
    so that the final model forward pass is: $h = W_0^{(l)} x + B_{\text{merged}}^{(l)}(\Theta) A_{\text{merged}}^{(l)}(\Theta) x$, ensuring zero increase in the parameters during forward passes.
    \item \textbf{Perplexity-Based Calibration:}
    For language models, the calibration dataset $\mathcal{D}_{\text{cal}} = \{(x_i^{(k)}, y_i^{(k)})\}_{i=1}^M$ consists of a tiny few-shot set of prompt-response pairs (e.g., 5--10 prompt sequences per task). The loss function $\ell_{\text{LLM}}$ is the autoregressive cross-entropy loss (token-level perplexity).
    
    The optimization problem is:
    \begin{equation}
    \min_{\Theta} \mathcal{L}_{\text{perplexity}}(\Theta) + \lambda_{\text{PAC}} \mathcal{R}_{\text{PAC}}(\Theta)
    \end{equation}
    where:
    \begin{equation}
    \mathcal{L}_{\text{perplexity}}(\Theta) = -\frac{1}{S \cdot K \cdot M} \sum_{s=1}^S \sum_{k=0}^{K-1} \sum_{i=1}^M \sum_{t=1}^{T_i} \ln P\left(y_{i,t}^{(k)} \mid y_{i,<t}^{(k)}, x_i^{(k)}; W_{\text{merged}}(\tilde{\Theta}_s)\right)
    \end{equation}
    This formulation guarantees that the merged LLM maintains high conversational fluency, avoids the destructive token-level interference ("hallucinations") typical of heuristic LoRA merging, and generalizes exceptionally well across all $K$ tasks under extreme calibration scarcity.
\end{enumerate}

\section{Derivation of General Non-Isotropic PAC-Bayesian Bounds with Diagonal Covariances and Fisher Information Regularization}
\label{app:non_isotropic_derivation}

In Section~\ref{subsec:pac_bayes_derivation}, we formulated the PAC-Bayesian prior and posterior distributions using isotropic Gaussian covariances ($\sigma_0^2 I$ and $\sigma^2 I$) for analytical simplicity. In this section, we lift this simplifying assumption and derive the general \emph{non-isotropic} PAC-Bayesian bound. We then mathematically prove how specifying a diagonal covariance matrix based on the empirical Fisher Information Matrix (FIM) of the network layers leads directly to a highly principled, layer-wise sensitive Consensus-Pulling penalty.

\subsection{General Diagonal Gaussian KL Divergence}
Let the learnable posterior distribution $Q$ and the prior distribution $P$ over the trajectory parameters $\tilde{\Theta} \in \mathbb{R}^{K(d+1)}$ be diagonal Gaussian distributions:
\begin{equation}
Q(\tilde{\Theta}) \sim \mathcal{N}\left(\Theta, \Sigma_{\text{post}}\right), \quad \Sigma_{\text{post}} = \mathrm{diag}\left(\sigma_{k, j}^2\right)
\end{equation}
\begin{equation}
P(\tilde{\Theta}) \sim \mathcal{N}\left(\Theta_{\text{uniform}}, \Sigma_{\text{prior}}\right), \quad \Sigma_{\text{prior}} = \mathrm{diag}\left(\sigma_{0, k, j}^2\right)
\end{equation}
where $k \in \{0, \dots, K-1\}$ indexing the task experts, and $j \in \{0, \dots, d\}$ indexing the polynomial trajectory coordinates.

The Kullback-Leibler divergence between these two multivariate Gaussians is solved as:
\begin{align}
D_{\text{KL}}(Q \parallel P) &= \int_{\mathbb{R}^{K(d+1)}} Q(\tilde{\Theta}) \ln \frac{Q(\tilde{\Theta})}{P(\tilde{\Theta})} d\tilde{\Theta} \\
&= \frac{1}{2} \left[ \ln \frac{|\Sigma_{\text{prior}}|}{|\Sigma_{\text{post}}|} - K(d+1) + \mathrm{Tr}\left(\Sigma_{\text{prior}}^{-1} \Sigma_{\text{post}}\right) + (\Theta - \Theta_{\text{uniform}})^T \Sigma_{\text{prior}}^{-1} (\Theta - \Theta_{\text{uniform}}) \right]
\end{align}
Since $\Sigma_{\text{prior}}$ and $\Sigma_{\text{post}}$ are diagonal matrices, the determinant and trace can be written as coordinate-wise sums:
\begin{equation}
\label{eq:diagonal_kl}
D_{\text{KL}}(Q \parallel P) = \frac{1}{2} \sum_{k=0}^{K-1} \sum_{j=0}^d \left( \frac{(\theta_{k,j} - \theta_{\text{uniform}, k, j})^2}{\sigma_{0, k, j}^2} + \frac{\sigma_{k, j}^2}{\sigma_{0, k, j}^2} - 1 - \ln \frac{\sigma_{k, j}^2}{\sigma_{0, k, j}^2} \right)
\end{equation}

By fixing the coordinate-wise posterior variances $\sigma_{k, j}^2$ during optimization, the terms involving $\sigma_{k, j}^2 / \sigma_{0, k, j}^2$ become constant. Minimizing McAllester's generalization bound with respect to the mean parameters $\Theta$ then reduces directly to minimizing the following coordinate-weighted regularizer:
\begin{equation}
\label{eq:weighted_regularizer}
\mathcal{R}_{\text{PAC}}(\Theta) = \sum_{k=0}^{K-1} \sum_{j=0}^d \frac{1}{\sigma_{0, k, j}^2} (\theta_{k,j} - \theta_{\text{uniform}, k, j})^2
\end{equation}

\subsection{Fisher Information Matrix (FIM) Regularization}
Equation~\eqref{eq:weighted_regularizer} proves that the strength of the Consensus-Pulling penalty for each trajectory coordinate is governed by its prior variance: coordinates with a small prior variance $\sigma_{0, k, j}^2$ are heavily penalized and pulled tightly toward the uniform baseline, while coordinates with a large prior variance are allowed to adapt more freely.

In deep neural networks, different layers and parameter coordinates exhibit highly heterogeneous sensitivities. To align our PAC-Bayesian bound with this sensitivity landscape, we define the prior variances using the empirical Fisher Information Matrix (FIM). Let $F_{k, j}$ represent the diagonal element of the empirical Fisher Information Matrix for the trajectory parameter $\theta_{k, j}$, evaluated over a small sample of the pre-trained base network $W_0$:
\begin{equation}
F_{k, j} = \mathbb{E}_{x \sim \mathcal{D}}\left[ \left( \frac{\partial \mathcal{L}_{\text{ce}}\left(f(x; W_{\text{merged}}(\Theta)); y\right)}{\partial \theta_{k, j}} \right)^2 \right]
\end{equation}
where the expectation is computed at the uniform baseline point $\Theta_{\text{uniform}}$. The Fisher diagonal $F_{k, j}$ represents the exact local sensitivity of the network loss to perturbations in coordinate $\theta_{k, j}$.

We specify the prior variances to be inversely proportional to their empirical Fisher sensitivities:
\begin{equation}
\sigma_{0, k, j}^2 = \frac{\tau^2}{F_{k, j} + \epsilon}
\end{equation}
where $\tau^2$ is a global prior scale parameter and $\epsilon > 0$ is a small damping coefficient to prevent division by zero in zero-gradient layers.

Substituting this prior variance back into our derived PAC-Bayesian regularizer in Equation~\eqref{eq:weighted_regularizer} yields:
\begin{equation}
\label{eq:fim_regularizer}
\mathcal{R}_{\text{PAC}}(\Theta) = \frac{1}{\tau^2} \sum_{k=0}^{K-1} \sum_{j=0}^d \left( F_{k, j} + \epsilon \right) \left( \theta_{k, j} - \theta_{\text{uniform}, k, j} \right)^2
\end{equation}

This is a highly significant theoretical result. Equation~\eqref{eq:fim_regularizer} proves that:
\begin{itemize}
    \item \textbf{Highly Sensitive Coordinates (Large $F_{k, j}$):} In layers that extract rich, representation-heavy shared features (e.g., intermediate Transformer layers), the Fisher sensitivity is high. The regularizer applies a heavy quadratic penalty, pulling these coordinates tightly back to the stable uniform ensembling consensus ($\theta_{k, j} \to \theta_{\text{uniform}, k, j}$). This mathematically suppresses destructive representation interference and preserves mode connectivity.
    \item \textbf{Insensitive Coordinates (Small $F_{k, j}$):} In layers with high task specificity or redundant representations (e.g., final classification heads or redundant residual projection layers), the Fisher sensitivity is low. The regularizer applies a very weak penalty, allowing these parameters to adjust freely and specialize on task-specific calibration features.
\end{itemize}
This FIM-regularized Consensus-Pulling penalty provides a rigorous, information-theoretic resolution to the simplifying assumption of isotropic variance, establishing a mathematically complete framework for heterogeneous parameter-space regularizations.
